\section{Zadanie 3}
\autor{Dominik Korsa}
Niech \(X\) oraz \(Y\) będą niezależnymi zmiennymi losowymi z rozkładu
jednostajnego na przedziale \([0, 1]\). Znajdź gęstość, dystrybuantę, wartość
oczekiwaną i wariancję zmiennych:

\begin{enumerate}[label=(\roman*)]
	\item\label{set08.task03.1} \(Z = X + Y\),
	\item\label{set08.task03.2} \(Z = X - Y\),
	\item\label{set08.task03.3} \(Z = |X - Y|\).
\end{enumerate}

\subsection*{Rozwiązanie \ref{set08.task03.1}}

Zmienne losowe \(X\) oraz \(Y\) mają rozkład opisany gęstością:
\[
	f_X(x) = f_Y(x)
	= \begin{cases}
		\frac{1}{1-0} = 1 & \text{w przedziale } x \in [0, 1] \\
		0                 & \text{wpp.}
	\end{cases}
\]
oraz dystrybuantą:
\[
	F_X(x) = F_Y(x)
	= \begin{cases}
		0                   & \text{gdy } x \leq 0              \\
		\frac{x-0}{1-0} = x & \text{w przedziale } x \in [0, 1] \\
		1                   & \text{gdy } x \geq 1
	\end{cases}.
\]

Skorzystamy z wzorów na splatanie z początku zestawu.
W rozwiązaniu podpunktu \ref{set08.task03.1} zacznę od policzenia gęstości \(f_{X+Y}\):
\begin{align*}
	f_{X+Y}(z)
	 & = \int_{-\infty}^{\infty} f_X(x) f_Y(z-x) \, dx            \\
	 & = \int_{0}^{1} f_X(x) f_Y(z-x) \, dx
	\tag{1}\label{set08.task03.1.tagA}                            \\
	 & = \begin{cases}
		     0
		      & \text{gdy \(z < 0\) lub \(z > 2\)} \\
		     \int_{0}^{z} f_X(x) f_Y(z-x) \, dx
		      & \text{gdy \(z \in [0, 1]\)}        \\
		     \int_{z-1}^{1} f_X(x) f_Y(z-x) \, dx
		      & \text{gdy \(z \in [1, 2]\)}        \\
	     \end{cases}
	\tag{2}\label{set08.task03.1.tagB}                            \\
	 & = \begin{cases}
		     0
		      & \text{gdy \(z < 0\) lub \(z > 2\)} \\
		     \int_{0}^{z} 1 \, dx
		      & \text{gdy \(z \in [0, 1]\)}        \\
		     \int_{z-1}^{1} 1 \, dx
		      & \text{gdy \(z \in [1, 2]\)}        \\
	     \end{cases}        \\
	 & = \boxed{\begin{cases}
			            0
			             & \text{gdy \(z < 0\) lub \(z > 2\)} \\
			            z
			             & \text{gdy \(z \in [0, 1]\)}        \\
			            2 - z
			             & \text{gdy \(z \in [1, 2]\)}        \\
		            \end{cases}} \\
\end{align*}

\eqref{set08.task03.1.tagA}, bo tylko gdy \(x \in (0, 1)\) zachodzi \(f_X(x) \neq 0\).

Dla \eqref{set08.task03.1.tagB} należy wyznaczyć dla jakich \(x \in [0, 1]\) zachodzi
również warunek \(z-x \in [0, 1]\), żeby \(f_X(z-x) \neq 0\).

Dystrybuanta \(F_{X+Y}(z)\) to z definicji:
\begin{align*}
	F_{X+Y}(z)
	 & = \int_{-\infty}^{z} f(z) \,dx                                 \\
	 & = \begin{cases}
		     0
		      & \text{gdy \(z < 0\)}        \\
		     \int_{0}^{z} x \, dx
		      & \text{gdy \(z \in [0, 1]\)} \\
		     \int_{0}^{1} x \, dx + \int_{1}^{z} 2-x \, dx
		      & \text{gdy \(z \in [1, 2]\)} \\
		     \int_{0}^{1} x \, dx + \int_{1}^{2} 2-x \, dx
		      & \text{gdy \(z > 2\)}        \\
	     \end{cases}                 \\
	 & = \begin{cases}
		     0
		      & \text{gdy \(z < 0\)}        \\
		     \frac{1}{2}z^2
		      & \text{gdy \(z \in [0, 1]\)} \\
		     \frac{1}{2} + \left(-\frac{1}{2}z^2 + 2z - \frac{3}{2}\right)
		      & \text{gdy \(z \in [1, 2]\)} \\
		     \frac{1}{2} + \frac{1}{2}
		      & \text{gdy \(z > 2\)}        \\
	     \end{cases} \\
	 & = \boxed{\begin{cases}
			            0
			             & \text{gdy \(z < 0\)}        \\
			            \frac{1}{2}z^2
			             & \text{gdy \(z \in [0, 1]\)} \\
			            1 - \frac{1}{2}(2-z)^2
			             & \text{gdy \(z \in [1, 2]\)} \\
			            1
			             & \text{gdy \(z > 2\)}        \\
		            \end{cases}}            \\
\end{align*}

Policzmy wartość oczekiwaną z liniowości:
\[
	E[X+Y] = E[X] + E[Y] = \frac{1-0}{2} + \frac{1-0}{2} = \boxed{1},
\]

oraz wariancję, dzięki niezależności zmiennych:
\[
	\variance{X+Y}
	= \variance{X} + \variance{Y}
	= 2 \cdot \frac{(1-0)^2}{12}
	= \boxed{\frac{1}{6}}
\]

\subsection*{Rozwiązanie \ref{set08.task03.2}}
Potraktujmy \(Z = X - Y\) jako sumę zmiennych losowych \(X\) oraz \(-Y\).
\(-Y\) możemy traktować jak zmienną losową z rozkładem jednostajnym na przedziale
\([-1, 0]\):
\[
	f_{-Y}(x)
	= \begin{cases}
		\frac{1}{0-(-1)} = 1 & \text{w przedziale } x \in [-1, 0] \\
		0                    & \text{wpp.}
	\end{cases}
\]
oraz:
\[
	F_{-Y}
	= \begin{cases}
		0                           & \text{gdy } x \leq -1              \\
		\frac{x-(-1)}{0-(-1)} = x+1 & \text{w przedziale } x \in [-1, 0] \\
		1                           & \text{gdy } x \geq 0
	\end{cases}.
\]

Policzmy gęstość \(f_{X-Y}\):
\begin{align*}
	f_{X-Y}(z)
	 & = \int_{-\infty}^{\infty} f_X(x) f_{-Y}(z-x) \, dx          \\
	 & = \begin{cases}
		     0
		      & \text{gdy \(z < -1\) lub \(z > 1\)} \\
		     \int_{0}^{1+z} f_X(x) f_Y(z-x) \, dx
		      & \text{gdy \(z \in [-1, 0]\)}        \\
		     \int_{z}^{1} f_X(x) f_Y(z-x) \, dx
		      & \text{gdy \(z \in [0, 1]\)}         \\
	     \end{cases}
	\tag{2}\label{set08.task03.2.tagA}                             \\
	 & = \begin{cases}
		     0
		      & \text{gdy \(z < -1\) lub \(z > 1\)} \\
		     \int_{0}^{1+z} 1 \, dx
		      & \text{gdy \(z \in [-1, 0]\)}        \\
		     \int_{z}^{1} 1 \, dx
		      & \text{gdy \(z \in [0, 1]\)}         \\
	     \end{cases}        \\
	 & = \boxed{\begin{cases}
			            0
			             & \text{gdy \(z < -1\) lub \(z > 1\)} \\
			            1+z
			             & \text{gdy \(z \in [-1, 0]\)}        \\
			            1-z
			             & \text{gdy \(z \in [0, 1]\)}         \\
		            \end{cases}} \\
\end{align*}

\eqref{set08.task03.2.tagA}: aby \(f_X(x) \neq 0\) oraz \(f_{-Y}(z-x) \neq 0\)
mamy warunki:
\begin{itemize}
	\item \(0 \leq x \leq 1\), oraz
	\item \(-1 \leq z-x \leq 0\), czyli \(z \leq x \leq z+1\).
\end{itemize}

Policzmy dysrybuantę:
\begin{align*}
	F_{X-Y}(z)
	 & = \int_{\infty}^z f(x) \, dx                      \\
	 & = \begin{cases}
		     0
		      & \text{gdy \(z < -1\)}        \\
		     \int_{-1}^{z} 1+x \, dx
		      & \text{gdy \(z \in [-1, 0]\)} \\
		     \int_{-1}^{0} 1+x \, dx + \int_{0}^{z} 1-x \, dx
		      & \text{gdy \(z \in [0, 1]\)}  \\
		     1
		      & \text{gdy \(z > 1\)}         \\
	     \end{cases} \\
	 & = \begin{cases}
		     0
		      & \text{gdy \(z < -1\)}        \\
		     \int_{0}^{z+1} x \, dx
		      & \text{gdy \(z \in [-1, 0]\)} \\
		     \int_{0}^{1} x \, dx - \int_{1}^{z+1} x \, dx
		      & \text{gdy \(z \in [0, 1]\)}  \\
		     1
		      & \text{gdy \(z > 1\)}         \\
	     \end{cases}    \\
	 & = \begin{cases}
		     0
		      & \text{gdy \(z < -1\)}        \\
		     \frac{1}{2}(z+1)^2
		      & \text{gdy \(z \in [-1, 0]\)} \\
		     1 - \frac{1}{2} (1-z)^2
		      & \text{gdy \(z \in [0, 1]\)}  \\
		     1
		      & \text{gdy \(z > 1\)}         \\
	     \end{cases}    \\
\end{align*}

Podpobnie jak w kroku \ref{set08.task03.1}, policzmy wartość oczekiwaną z liniowości:
\[
	E[X-Y] = E[X] - E[Y] = \frac{1-0}{2} - \frac{1-0}{2} = \boxed{0},
\]

oraz wariancję, dzięki niezależności zmiennych:
\[
	\variance{X+Y}
	= \variance{X} + \variance{-Y}
	= \frac{(1-0)^2}{12} + \frac{(0-1)^2}{12}
	= \boxed{\frac{1}{6}}
\]


\subsection*{Rozwiązanie \ref{set08.task03.3}}
\subsubsection*{Dystrybuanta \(F_Z(z)\)}
Rozkład łączny pary \((X, Y)\) jest jednostajny na kwadracie jednostkowym \([0, 1] \times [0, 1]\). Pole tego kwadratu wynosi 1.

Szukamy prawdopodobieństwa:
\[ F_Z(z) = P(Z \le z) = P(|X - Y| \le z) \]
Dla \(z \in [0, 1]\), nierówność \(|X - Y| \le z\) jest równoważna układowi:
\[ -z \le X - Y \le z \iff Y \le X + z \quad \text{oraz} \quad Y \ge X - z \]

Geometrycznie łatwiej jest obliczyć prawdopodobieństwo zdarzenia przeciwnego \(P(|X - Y| > z)\) i odjąć je od 1. Zdarzenie to odpowiada dwóm trójkątom w narożnikach kwadratu:
\begin{enumerate}
	\item Górny lewy róg (\(Y > X + z\)): trójkąt o przyprostokątnych \(1-z\). Pole wynosi \(\frac{1}{2}(1-z)^2\).
	\item Dolny prawy róg (\(Y < X - z\)): trójkąt o przyprostokątnych \(1-z\). Pole wynosi \(\frac{1}{2}(1-z)^2\).
\end{enumerate}

Suma pól obszarów "zakazanych" wynosi \((1-z)^2\). Zatem dystrybuanta dla \(z \in [0, 1]\) to:
\[ F_Z(z) = 1 - (1-z)^2 = 1 - (1 - 2z + z^2) = 2z - z^2 \]

Pełny wzór dystrybuanty:
\[
	F_Z(z) = \begin{cases}
		0        & \text{dla } z < 0         \\
		2z - z^2 & \text{dla } 0 \le z \le 1 \\
		1        & \text{dla } z > 1
	\end{cases}
\]

\subsubsection*{Gęstość \(f_Z(z)\)}
Gęstość otrzymujemy poprzez zróżniczkowanie dystrybuanty w przedziale \((0, 1)\):
\[ f_Z(z) = F_Z'(z) = (2z - z^2)' = 2 - 2z = 2(1-z) \]

Pełny wzór gęstości:
\[
	f_Z(z) = \begin{cases}
		2(1-z) & \text{dla } 0 < z < 1 \\
		0      & \text{w p.p.}
	\end{cases}
\]

\subsubsection*{Wartość oczekiwana \(E[Z]\)}
Obliczamy wartość oczekiwaną z definicji dla zmiennej ciągłej:
\[ E[Z] = \int_{-\infty}^{\infty} z \cdot f_Z(z) \, dz = \int_{0}^{1} z \cdot 2(1-z) \, dz \]
\[ = 2 \int_{0}^{1} (z - z^2) \, dz = 2 \left[ \frac{z^2}{2} - \frac{z^3}{3} \right]_0^1 \]
\[ = 2 \left( \frac{1}{2} - \frac{1}{3} \right) = 2 \cdot \frac{1}{6} = \frac{1}{3} \]

\subsubsection*{Wariancja \(Var(Z)\)}
Korzystamy ze wzoru \(Var(Z) = E[Z^2] - (E[Z])^2\). Najpierw obliczamy drugi moment:
\[ E[Z^2] = \int_{0}^{1} z^2 \cdot 2(1-z) \, dz = 2 \int_{0}^{1} (z^2 - z^3) \, dz \]
\[ = 2 \left[ \frac{z^3}{3} - \frac{z^4}{4} \right]_0^1 = 2 \left( \frac{1}{3} - \frac{1}{4} \right) = 2 \cdot \frac{1}{12} = \frac{1}{6} \]

Obliczamy wariancję:
\[ Var(Z) = \frac{1}{6} - \left(\frac{1}{3}\right)^2 = \frac{1}{6} - \frac{1}{9} = \frac{3}{18} - \frac{2}{18} = \frac{1}{18} \]


\begin{itemize}
	\item \textbf{Dystrybuanta:} \(F_Z(z) = 2z - z^2\) dla \(z \in [0,1]\)
	\item \textbf{Gęstość:} \(f_Z(z) = 2(1-z)\) dla \(z \in (0,1)\)
	\item \textbf{Wartość oczekiwana:} \(E[Z] = \frac{1}{3}\)
	\item \textbf{Wariancja:} \(Var(Z) = \frac{1}{18}\)
\end{itemize}